\documentclass[11pt,a4paper]{article}
\usepackage[margin=25mm,headheight=14pt]{geometry}
\usepackage{fontspec}
\setmainfont{lmroman10-regular.otf}[BoldFont=lmroman10-bold.otf,ItalicFont=lmroman10-italic.otf,BoldItalicFont=lmroman10-bolditalic.otf]
\setsansfont{lmsans10-regular.otf}
\setmonofont{lmmono10-regular.otf}[Scale=MatchLowercase]
\newfontfamily\authorsymbol{Noto Sans Symbols}[Scale=MatchUppercase]
\usepackage{amsmath,amssymb,amsthm,mathtools,booktabs}
\usepackage{microtype}
\usepackage[hidelinks,unicode]{hyperref}
\usepackage{fancyhdr}
\pagestyle{fancy}\fancyhf{}\fancyhead[L]{\small\itshape Affine cycles and minimal retained state}\fancyhead[R]{\small\thepage}\renewcommand{\headrulewidth}{0pt}
\setlength{\parindent}{1.2em}\setlength{\parskip}{0pt}\setlength{\emergencystretch}{2em}
\newtheorem{theorem}{Theorem}\newtheorem{proposition}[theorem]{Proposition}\newtheorem{lemma}[theorem]{Lemma}\newtheorem{corollary}[theorem]{Corollary}
\theoremstyle{definition}\newtheorem{definition}[theorem]{Definition}\newtheorem{example}[theorem]{Example}\newtheorem{remark}[theorem]{Remark}
\DeclareMathOperator{\rank}{rank}\DeclareMathOperator{\im}{im}\DeclareMathOperator{\spanop}{span}\DeclareMathOperator{\diag}{diag}
\newcommand{\R}{\mathbb R}\newcommand{\Q}{\mathbb Q}\newcommand{\Z}{\mathbb Z}\newcommand{\N}{\mathbb N}\newcommand{\id}{\mathrm{id}}
\hypersetup{pdftitle={Affine Cycles: Exact Cut Defects and Minimal Retained State},pdfauthor={lmtlssss🜐},pdfsubject={Mathematics: affine composition and minimal linear summaries},pdfkeywords={affine maps, cyclic composition, accumulated costs, minimal linear state, exact arithmetic}}
\title{\LARGE Affine Cycles:\\Exact Cut Defects and Minimal Retained State}
\author{lmtlssss{\authorsymbol 🜐}}
\date{3 November 2026\\[3pt]\small\href{https://doi.org/10.5281/zenodo.23931722}{10.5281/zenodo.23931722}}
\begin{document}\maketitle\thispagestyle{plain}

\begin{abstract}
An ordered cycle of affine maps carries a scalar cost at each step. We
compare each cyclic rotation at the initial state transported to its cut.
For affine costs, the complete family of cut defects is exactly
$R((M-I)x+b)$, where $x\mapsto Mx+b$ is the full cycle and the rows of
$R$ are the accumulated prefix-cost coefficients. Thus the minimum
linear summary of the initial state has dimension $\rank(R(M-I))$,
even when the decoder is arbitrary. The affine system
$R(M-I)x=-Rb$ classifies all cut-invariant states. It can have solutions
when the full cycle has no fixed point. We also give the minimum summary
for all transported cut totals, exact interval-substitution formulas,
and polynomial-cost boundary identities. For polynomial defects,
a finite matrix of derivative coefficients determines the minimum
linear dimension. A rational three-dimensional example has a
one-dimensional defect summary and a plane of cut-invariant states
without a fixed point. Exact coefficient certificates and a
source-derived mathematical figure accompany the proofs.
\end{abstract}
\noindent\textit{Keywords.} Affine maps; cyclic composition; accumulated
costs; minimal linear state; exact arithmetic; polynomial invariance.

\section{An ordered cycle and its transported cuts}
Fix a dimension $d\ge1$ and a word of length $n\ge1$. For
$0\le i<n$, let
\[
 f_i(x)=A_i x+b_i,\qquad c_i(x)=a_i^{\mathsf T}x+\alpha_i,
 \qquad x\in\R^d.
\]
The cost $c_i$ is evaluated at the state \emph{before} applying $f_i$.
No sign condition is imposed on the scalar costs. Define the prefix
state and prefix cost by
\begin{equation}\label{eq:prefix}
 \begin{aligned}
 P_0(x)&=x,&P_k&=f_{k-1}\circ\cdots\circ f_0,\\
 C_0(x)&=0,&C_k(x)&=\sum_{i=0}^{k-1}c_i(P_i(x)).
 \end{aligned}
\end{equation}
Write
\[
 P_k(x)=M_kx+p_k,\qquad C_k(x)=r_k^{\mathsf T}x+\sigma_k,
 \qquad F=P_n,
\]
and abbreviate $M=M_n$, $b=p_n$, $C=C_n$, and
$N=M-I$. The full-cycle displacement is
\begin{equation}\label{eq:displacement}
 \delta(x)=F(x)-x=Nx+b.
\end{equation}
Every prefix is defined by composition, so these expressions remain
valid when a local or prefix matrix is singular.

A cut at $k$, with $0\le k<n$, puts the steps in the order
$k,k+1,\ldots,n-1,0,\ldots,k-1$. Let $C^{[k]}(y)$ denote the total cost
of that rotated word when its initial state is $y$.
The \emph{transported cut total} and the \emph{cut defect} are
\begin{equation}\label{eq:defectdef}
 T_k(x)=C^{[k]}(P_k(x)),\qquad D_k(x)=T_k(x)-C(x).
\end{equation}
In particular, $T_0=C$ and $D_0=0$. The comparison in
\eqref{eq:defectdef} uses the state reached at the cut. A comparison
that resets every rotated word to the same initial state defines a
different family of functions.

We call $x$ \emph{cut-invariant} when $D_k(x)=0$ for all $k$.
This concerns the accumulated costs. A fixed point satisfies the
additional state equation $F(x)=x$.

\section{The exact cut-defect identity}
The boundary identity uses only the addition of local costs. We first
state it in a form that also applies beyond affine costs.

\begin{lemma}[Boundary identity]\label{lem:boundary}
For arbitrary state maps and scalar local costs for which the
compositions are defined,
\begin{equation}\label{eq:boundary}
 \boxed{D_k(x)=C_k(F(x))-C_k(x),\qquad 0\le k<n.}
\end{equation}
\end{lemma}
\begin{proof}
Let $S_k=f_{n-1}\circ\cdots\circ f_k$ be the suffix state map.
Let $B_k(y)$ be the cost of that suffix, started at $y$.
The original word consists of the prefix followed by the suffix, hence
\[
 F=S_k\circ P_k,\qquad C(x)=C_k(x)+B_k(P_k(x)).
\]
The rotated word consists of the same suffix followed by the same
prefix. Starting it at $P_k(x)$ therefore gives
\[
 T_k(x)=B_k(P_k(x))+C_k(S_k(P_k(x)))
       =B_k(P_k(x))+C_k(F(x)).
\]
Subtracting the preceding expression for $C(x)$ proves
\eqref{eq:boundary}. This argument uses composition and addition and
requires no inverse map.
\end{proof}

Stack the nontrivial prefix rows into
\begin{equation}\label{eq:RK}
 R=\begin{bmatrix}r_1^{\mathsf T}\\\vdots\\r_{n-1}^{\mathsf T}\end{bmatrix},
 \qquad K=RN,\qquad h=Rb,
 \qquad D(x)=(D_1(x),\ldots,D_{n-1}(x))^{\mathsf T}.
\end{equation}
For $n=1$, these are empty-row matrices and an empty vector, with rank
zero.

\begin{theorem}[Cut defects and cut-invariant states]\label{thm:cut}
For affine state maps and affine local costs,
\begin{equation}\label{eq:main}
 \boxed{D(x)=R\delta(x)=Kx+h.}
\end{equation}
The cut-invariant set is nonempty exactly when
\begin{equation}\label{eq:consistency}
 h\in\im K,
 \quad\text{equivalently}\quad
 \rank[\,K\mid h\,]=\rank K.
\end{equation}
When it is nonempty, it is the affine space
\begin{equation}\label{eq:zero}
 x_*+\ker K,\qquad Kx_*=-h,
\end{equation}
of dimension $d-\rank K$. Every fixed point is cut-invariant.
\end{theorem}
\begin{proof}
By Lemma~\ref{lem:boundary} and the affine form of $C_k$,
\[
 D_k(x)=r_k^{\mathsf T}(Mx+b)+\sigma_k
              -r_k^{\mathsf T}x-\sigma_k
       =r_k^{\mathsf T}(Nx+b).
\]
Stacking these identities proves \eqref{eq:main}. Cut invariance is
therefore exactly the linear system $Kx=-h$. Its solvability,
solution space and dimension are given by \eqref{eq:consistency},
\eqref{eq:zero}, and rank--nullity. A fixed point has
$\delta(x)=0$, so \eqref{eq:main} makes every defect zero.
\end{proof}

The entire prefix coefficient family is applied only to the attainable
affine displacement set $b+\im N$. Directions outside this set do not
increase the dimension needed to distinguish the defects.
The offset $h$ affects the zero-defect equation, but its presence does
not add an unknown coordinate to the initial state.

\begin{corollary}[When cost invariance forces state closure]\label{cor:closure}
Suppose the full cycle has a fixed point $x_0$. The cut-invariant set
coincides with the fixed-point set exactly when
\begin{equation}\label{eq:observability}
 \rank(RN)=\rank N.
\end{equation}
In particular, this holds if $R$ is injective on $\im N$.
\end{corollary}
\begin{proof}
Since $Nx_0=-b$, the fixed-point set is $x_0+\ker N$ and
Theorem~\ref{thm:cut} makes the cut-invariant set $x_0+\ker(RN)$.
The inclusion $\ker N\subseteq\ker(RN)$ is automatic. The two kernels
are equal exactly when their dimensions are equal, which is
\eqref{eq:observability}. This is also the condition that the
restriction of $R$ to $\im N$ be injective.
\end{proof}

\section{Minimum linear state under an explicit recovery contract}
We measure the number of real coordinates retained by a linear summary.
The map family and its coefficients are fixed and known to the decoder.

\begin{definition}[Linear summary]\label{def:summary}
For a function $H:\R^d\to\R^m$, a linear summary of dimension $q$ is
a linear map $S:\R^d\to\R^q$ and a function
$g:\R^q\to\R^m$ such that $H(x)=g(Sx)$ for every $x\in\R^d$.
The decoder $g$ may be arbitrary. Dimension zero denotes a constant
summary.
\end{definition}

\begin{lemma}[Affine recovery rank]\label{lem:rank}
For $H(x)=Lx+a$, the minimum linear-summary dimension is $\rank L$.
\end{lemma}
\begin{proof}
If $v\in\ker S$, then $S(x+v)=Sx$ for every $x$; the recovery
contract gives $H(x+v)=H(x)$, hence $Lv=0$. Thus
$\ker S\subseteq\ker L$ and $q\ge\rank S\ge\rank L$.
Conversely, let the rows of $S$ be a basis of the row space of $L$.
There is a matrix $B$ with $L=BS$, and the decoder
$g(z)=Bz+a$ recovers $H$. This uses $\rank L$ coordinates.
\end{proof}

\begin{theorem}[Minimum retained state]\label{thm:minimal}
Under Definition~\ref{def:summary}, the exact minimum dimensions are
\begin{align}
 q_{\mathrm{defect}}&=\rank\bigl(R(M-I)\bigr),\label{eq:minD}\\
 q_{\mathrm{total}}&=\rank\begin{bmatrix}r_n^{\mathsf T}\\R(M-I)\end{bmatrix}.
 \label{eq:minT}
\end{align}
The first formula recovers all cut defects. The second recovers all
transported cut totals $(T_0,\ldots,T_{n-1})$. A basis of the respective
row space gives an attaining summary in each case.
\end{theorem}
\begin{proof}
Equation~\eqref{eq:main} and Lemma~\ref{lem:rank} give
\eqref{eq:minD}. The vector of transported totals has linear part
\[
 L_T=\begin{bmatrix}
 r_n^{\mathsf T}\\
 r_n^{\mathsf T}+r_1^{\mathsf T}N\\
 \vdots\\
 r_n^{\mathsf T}+r_{n-1}^{\mathsf T}N
 \end{bmatrix}.
\]
Subtract its first row from every other row. These elementary row
operations preserve rank and yield the matrix in \eqref{eq:minT}.
Lemma~\ref{lem:rank} then proves both the lower bound and attainability.
\end{proof}

The retained state is a linear image of the initial state, and recovery
holds over all real initial states. The dimension counts real linear
coordinates. Arbitrary nonlinear encodings, bit storage and recovery
on a finite input set define separate contracts. Recovery of $Ru$
for every unconstrained displacement $u\in\R^d$ has minimum
dimension $\rank R$. Recovery on the attainable displacements of
the given cycle has the generally smaller rank in \eqref{eq:minD}.
This explicit recovery contract follows the distinction used for
minimal linear state in the preceding cyclic-shift note~\cite{cyclic}.

\section{Chronological composition and exact substitutions}
An affine word is fully represented by a tuple
$(A,b,r,\sigma)$ with state map $x\mapsto Ax+b$ and cost
$x\mapsto r^{\mathsf T}x+\sigma$.
The chronological product means that the left tuple acts first.

\begin{proposition}[Affine state and cost composition]\label{prop:compose}
The chronological product is
\begin{equation}\label{eq:product}
 (A,b,r,\sigma)\star(B,c,s,\tau)
 =\bigl(BA,\ Bb+c,\ r+A^{\mathsf T}s,\
             \sigma+\tau+s^{\mathsf T}b\bigr).
\end{equation}
It is associative, with identity $(I,0,0,0)$.
\end{proposition}
\begin{proof}
The resulting state is $B(Ax+b)+c$. The resulting cost is
$r^{\mathsf T}x+\sigma+s^{\mathsf T}(Ax+b)+\tau$, which gives
\eqref{eq:product}. For any two represented pairs $(f,c)$ and $(g,e)$,
the product represents $(g\circ f,c+e\circ f)$. Applying a third pair
$(h,a)$ in either grouping gives the state $h\circ g\circ f$ and the
cost $c+e\circ f+a\circ g\circ f$. Thus the product is associative.
The identity follows from the same representation.
\end{proof}

For homogeneous coordinates $\widetilde x=(x,1)^{\mathsf T}$, put
\[
 H_f=\begin{bmatrix}A&b\\0&1\end{bmatrix},
 \qquad \lambda_c=[\,r^{\mathsf T}\ \sigma\,].
\]
Then \eqref{eq:product} is simply
$(H_f,\lambda_c)\star(H_g,\lambda_e)
=(H_gH_f,\lambda_c+\lambda_eH_f)$.
This representation also makes the chronological order explicit.

\begin{proposition}[One interval substitution]\label{prop:substitution}
Let a word be the prefix $P$, then an interval $V$, then the suffix $S$.
Use $P,V,S$ also for their state maps, and write $C_P,C_V,C_S$ for
their costs. Replace $V,C_V$ by $V',C_{V'}$. With $y=P(x)$, the exact
change in total cost is
\begin{equation}\label{eq:substitution}
 C'(x)-C(x)=C_{V'}(y)-C_V(y)+C_S(V'(y))-C_S(V(y)).
\end{equation}
If $C_S(z)=s^{\mathsf T}z+\tau$, its last two terms reduce to
$s^{\mathsf T}(V'(y)-V(y))$.
If the suffix state map has linear part $A_S$, the final state changes by
$A_S(V'(y)-V(y))$.
\end{proposition}
\begin{proof}
The original total cost is $C_P(x)+C_V(y)+C_S(V(y))$; its replacement
is $C_P(x)+C_{V'}(y)+C_S(V'(y))$. Their difference gives
\eqref{eq:substitution}. Subtraction cancels the constant part of an
affine suffix cost or suffix state map, yielding the remaining formulas.
\end{proof}

Associativity permits a balanced binary tree of ordered interval
summaries. A tree with $n$ leaves stores $2n-1$ summaries; its root is
the full word. A leaf replacement requires recomputation only on its
ancestor path. A recursively balanced tree has height
$\lceil\log_2 n\rceil$, so one replacement changes at most that many
internal products. For $k$ replaced leaves the number of distinct
changed internal nodes is at most
$\min\{n-1,k\lceil\log_2 n\rceil\}$.
These are algebraic product counts; the cost of rational arithmetic
also depends on coefficient bit length. Ordered associative prefix
products and their tree organization are established constructions;
see Blelloch~\cite{blelloch}. The formulas above specify the state and
cost algebra to which they apply.

\section{Polynomial costs and the exact derivative criterion}
Lemma~\ref{lem:boundary} remains valid for arbitrary local costs.
For affine state maps and polynomial local costs of degree at most
$p$, every prefix cost and every cut defect has degree at most $p$.

For quadratic costs, represent a word by $(A,b,Q,u,v)$, where
$Q=Q^{\mathsf T}$ and
\[
 f(x)=Ax+b,\qquad c(x)=x^{\mathsf T}Qx+2u^{\mathsf T}x+v.
\]
Composing this word first and $(B,c,K,w,z)$ second gives
\begin{equation}\label{eq:quadraticproduct}
 \begin{aligned}
 A_{\rm new}&=BA,& b_{\rm new}&=Bb+c,\\
 Q_{\rm new}&=Q+A^{\mathsf T}KA,
 &u_{\rm new}&=u+A^{\mathsf T}(Kb+w),\\
 v_{\rm new}&=v+b^{\mathsf T}Kb+2w^{\mathsf T}b+z.
 \end{aligned}
\end{equation}
These identities follow by expanding the second cost at $Ax+b$.
Their associativity follows from the same state-and-cost
representation as Proposition~\ref{prop:compose}.

If $C_k(x)=x^{\mathsf T}Q_kx+2u_k^{\mathsf T}x+v_k$, then the boundary
identity yields
\begin{equation}\label{eq:quadraticboundary}
 D_k(x)=2(Q_kx+u_k)^{\mathsf T}\delta(x)
                    +\delta(x)^{\mathsf T}Q_k\delta(x).
\end{equation}
Equivalently, $D_k(x)=x^{\mathsf T}H_kx+2g_k^{\mathsf T}x+\beta_k$, with
\begin{equation}\label{eq:quadraticcoeff}
 \begin{aligned}
 H_k&=M^{\mathsf T}Q_kM-Q_k,\\
 g_k&=M^{\mathsf T}(Q_kb+u_k)-u_k,\\
 \beta_k&=b^{\mathsf T}Q_kb+2u_k^{\mathsf T}b.
 \end{aligned}
\end{equation}
Unlike an affine defect, a quadratic defect can depend on the initial
state in addition to its cycle displacement. The following theorem
still gives an exact finite rank criterion.

\begin{theorem}[Polynomial minimum linear dimension]\label{thm:polynomial}
Let $H=(H_1,\ldots,H_m):\R^d\to\R^m$ be a polynomial map.
Form a matrix $\mathcal G_H$ whose rows are indexed by pairs
$(i,\alpha)$ and whose $j$th entry in that row is the coefficient
of $x^\alpha$ in $\partial H_i/\partial x_j$. Include all monomials
that occur in these derivatives. Then
\begin{equation}\label{eq:polyrank}
 \boxed{q_{\min}(H)=\rank\mathcal G_H.}
\end{equation}
For quadratic cut defects with coefficients \eqref{eq:quadraticcoeff},
this becomes
\begin{equation}\label{eq:quadrank}
 q_{\mathrm{defect}}=
 \rank\begin{bmatrix}H_1\\g_1^{\mathsf T}\\\vdots\\H_{n-1}\\g_{n-1}^{\mathsf T}\end{bmatrix}.
\end{equation}
\end{theorem}
\begin{proof}
For $v\in\R^d$, let $\partial_v H_i=\sum_jv_j\partial H_i/\partial x_j$.
The definition of $\mathcal G_H$ gives
\[
 U:=\ker\mathcal G_H
   =\{v:\partial_vH_i\equiv0\text{ for every }i\}.
\]
If $H(x)=g(Sx)$ and $v\in\ker S$, then
$H_i(x+tv)=H_i(x)$ for all $x,t$. Differentiate this polynomial
identity at $t=0$. It gives $v\in U$, so
$\ker S\subseteq U$ and $q\ge d-\dim U=\rank\mathcal G_H$.
This step does not differentiate the arbitrary decoder.

Conversely, if $v\in U$, then the polynomial
$t\mapsto H_i(x+tv)$ has zero derivative for all real $t$, and is
constant. Hence $H$ is constant on each coset of $U$. Choose a
surjective linear map $S:\R^d\to\R^{d-\dim U}$ with kernel $U$.
Define $g(Sx)=H(x)$. Constancy on cosets makes $g$ well-defined and
proves attainability. For a quadratic component,
$\partial_vD_k(x)=2x^{\mathsf T}H_kv+2g_k^{\mathsf T}v$.
It vanishes identically exactly when $H_kv=0$ and
$g_k^{\mathsf T}v=0$. Stacking these conditions proves
\eqref{eq:quadrank}.
\end{proof}

\section{A cut-invariant plane without a fixed point}
We now give all local maps and costs for the principal example.
Coordinates $(x,y,z)$ denote the input to the indicated local step.
\begin{equation}\label{eq:examplemaps}
 \begin{array}{c|c|c}
 i&f_i(x,y,z)&c_i(x,y,z)\\\hline
 0&(x+y,\ y,\ z)&x+y\\
 1&(x,\ y+1,\ z+x)&x+z\\
 2&(x,\ y,\ z-y)&3x-2z\\
 3&(2x-2y+2,\ y,\ z-x+y)&-2x+4y+3z
 \end{array}
\end{equation}
The prefix state maps are
\[
 \begin{aligned}
 P_1(x,y,z)&=(x+y,y,z),\\
 P_2(x,y,z)&=(x+y,y+1,z+x+y),\\
 P_3(x,y,z)&=(x+y,y+1,z+x-1),\\
 F(x,y,z)&=(2x,y+1,z).
 \end{aligned}
\]
Direct summation gives
\begin{equation}\label{eq:examplecosts}
 \begin{aligned}
 C_1&=x+y,& C_2&=2x+2y+z,\\
 C_3&=3x+3y-z,& C&=4x+5y+2z+1.
 \end{aligned}
\end{equation}
Thus
\begin{equation}\label{eq:examplematrices}
 R=\begin{bmatrix}1&1&0\\2&2&1\\3&3&-1\end{bmatrix},\quad
 N=\diag(1,0,0),\quad b=\begin{bmatrix}0\\1\\0\end{bmatrix},\quad
 K=\begin{bmatrix}1&0&0\\2&0&0\\3&0&0\end{bmatrix}.
\end{equation}
The prefix row rank is two, whereas $\rank K=1$ and
$h=(1,2,3)^{\mathsf T}$. Consequently,
\begin{equation}\label{eq:exampledefects}
 \boxed{D(x,y,z)=(x+1)(1,2,3)^{\mathsf T}.}
\end{equation}
The exact linear summary is $S(x,y,z)=x$; its decoder multiplies
$x+1$ by $(1,2,3)^{\mathsf T}$. Cut-invariant states form the entire
plane $x=-1$. The equation $F(x,y,z)=(x,y,z)$ requires $y+1=y$, so
this cycle has no fixed point. Along the cut-invariant plane, its
nonzero displacement is $(-1,1,0)$, which lies in $\ker R$.

For all transported cut totals, Theorem~\ref{thm:minimal} gives rank two:
\[
 \rank\begin{bmatrix}4&5&2\\1&0&0\\2&0&0\\3&0&0\end{bmatrix}=2.
\]
An attaining summary is $(x,5y+2z)$. The first coordinate determines
the defects, while both coordinates determine $C$ in
\eqref{eq:examplecosts} and hence every transported cut total.

For the rational initial state $(2,-1,3)$, a direct pass through the
four original local steps gives costs $1,4,-5,10$, with total $10$.
Running the rotated words at their transported states gives
\[
 (T_0,T_1,T_2,T_3)=(10,13,16,19).
\]
The defects $(3,6,9)$ agree with \eqref{eq:exampledefects}.
At $(-1,0,0)$, all four transported totals equal $-3$.

The accompanying figure enumerates all $343$ states in
$\{-3,-2,-1,0,1,2,3\}^3$. Their cut defects take exactly seven
values, indexed by $t=x+1\in\{-2,-1,0,1,2,3,4\}$.
Each value has exactly $49$ preimages in this grid. The $49$ gold
states form the zero-defect fibre $x=-1$. The figure's scalar count
$q_{\mathrm{defect}}=1$ follows from the theorem for all real states;
the finite grid displays its fibres.

\section{Consistency, singular prefixes and polynomial separation}
\begin{example}[Affine inconsistency]\label{ex:inconsistent}
Keep the maps in \eqref{eq:examplemaps} and replace
$c_2(x,y,z)$ by $3x+y-2z$. The prefix costs become
\[
 C_1=x+y,\quad C_2=2x+2y+z,\quad C_3=3x+4y-z+1.
\]
The new defects are
$(x+1,2x+2,3x+4)^{\mathsf T}$. Their linear coefficient matrix
still has rank one, but adjoining the constant column gives rank two.
The first defect forces $x=-1$, at which the third defect is one.
Thus the cut-invariant set is empty. The full cost changes by
$y+1$, as Proposition~\ref{prop:substitution} also gives directly.
\end{example}

\begin{example}[A singular prefix]\label{ex:singular}
On $\R^2$, take
\[
 f_0(x,y)=(x,0),\quad c_0(x,y)=x+y,\qquad
 f_1(x,y)=(0,y+1),\quad c_1(x,y)=2x+y.
\]
Here $P_1$ is singular, $F(x,y)=(0,1)$ and $C(x,y)=3x+y$.
The rotated word started at $P_1(x,y)=(x,0)$ has cost $2x+1$.
Therefore $D_1(x,y)=1-x-y$, exactly as
$r_1^{\mathsf T}\delta(x,y)$ predicts. The cut-invariant line
$x+y=1$ contains the single fixed point $(0,1)$.
\end{example}

\begin{example}[Quadratic defects distinguish two rank questions]\label{ex:quadratic}
On $\R^2$, let $f_0$ be the identity with cost $c_0(x,y)=xy$;
let the second cost be zero. If $f_1(x,y)=(x+1,y)$, then
$D_1(x,y)=y$, with minimum linear-summary dimension one.
If instead $f_1(x,y)=(2x,y)$, then $D_1(x,y)=xy$.
Its displacement linear part has rank one, while its derivative
coefficient matrix has rank two: $\partial_x(xy)=y$ and
$\partial_y(xy)=x$. Theorem~\ref{thm:polynomial} proves that two
linear coordinates are necessary and sufficient. Thus the affine
rank formula has exactly the stated affine-cost scope.
\end{example}

\section{Exact certificates and relation to established methods}
The verifier uses rational arithmetic. One construction composes state
and cost coefficients by \eqref{eq:product} and
\eqref{eq:quadraticproduct}. A second, independent construction expands
local trajectories as multivariate polynomials and executes each
rotation at its symbolically transported state. Coefficient equality
establishes each certified identity for every initial state of the
specified word.

The evidence covers five named examples and $64$ deterministically
generated rational words, of dimensions one through three and lengths
one through six. These $69$ words supply $165$ nontrivial exact cut
identities and their zero-cut normalizations. Exact row reduction
checks the ranks and consistency conditions. Direct trajectory
enumeration checks all $343$ figure states and their four transported
cut totals. The source includes the generator, seed and coefficient
construction. The certificates establish these examples; the proofs
establish the general theorems.

Affine composition, rank--nullity, polynomial differentiation and
associative prefix products are established ingredients. This note
combines them for the transported-cut comparison: the boundary defect,
its restricted displacement rank, the consistency classification,
and the separate minimum dimensions for defects and totals. The
polynomial derivative criterion extends the same linear-summary
contract. The preceding cyclic-shift note specifies a related recovery
contract~\cite{cyclic}; Blelloch gives the broader associative-prefix
framework~\cite{blelloch}.

\begin{thebibliography}{9}
\bibitem{cyclic} lmtlssss{\authorsymbol 🜐},
\emph{Cyclic-Shift Energies: Minimal State and Cyclotomic Gaps},
version 1.0.0, Zenodo (2026).
\href{https://doi.org/10.5281/zenodo.23897127}{doi:10.5281/zenodo.23897127}.
\bibitem{blelloch} G. E. Blelloch,
\emph{Prefix Sums and Their Applications}, Technical Report
CMU-CS-90-190, School of Computer Science, Carnegie Mellon University
(November 1990).
\href{https://www.cs.cmu.edu/afs/cs.cmu.edu/project/scandal/public/papers/CMU-CS-90-190.html}{Official report record}.
\end{thebibliography}
\end{document}
